\documentclass[11pt]{article}
\usepackage[margin=1in]{geometry}
\usepackage{amsmath}
\begin{document}
\thispagestyle{empty}
\noindent Editorial Office\\
IEEE Transactions on Network and Service Management

\bigskip
\noindent Dear Editor,

\bigskip
\noindent Please consider the enclosed manuscript, ``When Does Recalibration Help Cellular Throughput Commitments? Accounting for Cost, Missingness, and Abstention,'' formatted in two-column \texttt{IEEEtran} journal style for consideration in \emph{IEEE Transactions on Network and Service Management}.

\bigskip
\noindent Centered on a practical empirical accounting methodology, reporting checklist, and reusable open-source CLI tool (\texttt{audit\_policy.py}), the manuscript leads with our descriptive finding across 16 retrospective physical campaign-dates and 26 held-out archive dates under frozen predictors and fixed policies: mobile conditional upper violation risk spans $9\%$--$68\%$ above nominal $\alpha=0.10$ while supported measured throughput remains around $1\%$--$3\%$, aggregate contrasts depend heavily on offer-count pooling weights ($-23.41$ vs.\ $+0.37$~pp), and sparse finite-sample bounds become vacuous on $36/40$ evaluated cells. Separating equal-cluster and pooled-ratio relative TOST equivalence at a $\pm 50\%$ baseline margin shows that naive date-level equivalence largely fails or collapses under three-campaign clustering sensitivity ($G_{\text{eff}}=3.84/16$ retrospective MARUR; $6.73/26$ held-out cap-40 PACER, $p=0.0071 \to 0.1153$), leaving utility differences largely inconclusive.

\bigskip
\noindent All evidence is retrospective from public archives, with no live field validation, and it does not cover other devices, operators, measurement tools, or closed-loop deployment. Repository source code is released under the MIT License (\texttt{LICENSE}), and third-party public dataset terms are documented separately without redistribution.

\bigskip
\noindent Sincerely,\\
Md.\ Moazzem Hossain Majumder
\end{document}
